\documentclass[conference]{IEEEtran}
\usepackage{cite}
\usepackage{amsmath,amssymb,amsfonts}
\usepackage{graphicx}
\usepackage{textcomp}
\usepackage{xcolor}
\usepackage{booktabs}
\usepackage{hyperref}
\usepackage{algorithm}
\usepackage{algorithmic}
\usepackage{multirow}

\begin{document}

\title{Proactive Kubernetes Autoscaling Using Patch-Based Transformer Models for Latency-Sensitive Microservices}

\author{\IEEEauthorblockN{Dhruv\IEEEauthorrefmark{1} and Sougata\IEEEauthorrefmark{2}}
\IEEEauthorblockA{\IEEEauthorrefmark{1}\IEEEauthorrefmark{2}\textit{Department of Computer Science} \\
\textit{Indian Institute of Technology Delhi}\\
New Delhi, India \\
dhruv@iitd.ac.in}
}

\maketitle

% ==============================================================================
% ABSTRACT
% ==============================================================================
\begin{abstract}
The Kubernetes Horizontal Pod Autoscaler (HPA) follows a reactive control loop that scales resources only after threshold violations are observed, introducing a systemic lag of 60--90 seconds during sudden traffic surges. This delay is particularly damaging in latency-sensitive microservice deployments, where even brief periods of under-provisioning degrade user experience and violate service-level objectives. We present a proactive autoscaling framework that replaces the reactive feedback loop with a predictive controller built on PatchTST, a patch-based Transformer architecture for time-series forecasting. Our evaluation spans two heterogeneous microservice benchmarks---the Google Online Boutique (10 services, Go/Python) and Weaveworks Sock Shop (13+ services, Java/Node.js/Go)---deployed on Google Kubernetes Engine clusters with Prometheus-based observability at 2--15 second resolution. Over 117 hours of continuous traffic simulation (28,108 temporally aligned observations for Online Boutique and 7,147 for Sock Shop), we demonstrate that PatchTST achieves an RMSE of 0.0179 on Online Boutique and 0.0917 on Sock Shop, representing a 97.7\% and 85.5\% reduction in mean absolute error compared to the reactive HPA baseline, respectively. The proactive controller detects 98.2\% of traffic spikes on Online Boutique and 89.5\% on Sock Shop, while maintaining false positive rates below 5\%---compared to HPA's 52\% recall and 43--48\% resource waste under identical conditions. A zero-shot transfer learning experiment, in which a model trained exclusively on Online Boutique data is evaluated on Sock Shop without fine-tuning, yields a 69\% improvement over HPA, confirming that the learned representations generalize across architecturally distinct microservice platforms. These results establish patch-based Transformer forecasting as a robust and generalizable alternative to threshold-driven autoscaling in production Kubernetes environments.
\end{abstract}

\begin{IEEEkeywords}
Kubernetes, Autoscaling, Transformer, PatchTST, Time-Series Forecasting, Cloud Engineering, Microservices, Transfer Learning
\end{IEEEkeywords}

% ==============================================================================
\section{Introduction}
% ==============================================================================

Microservice architectures have become the dominant paradigm for cloud-native applications, offering modularity, independent scalability, and fault isolation~\cite{dragoni2017microservices}. Kubernetes has emerged as the de facto orchestration platform, with its Horizontal Pod Autoscaler (HPA) providing automated scaling based on observed resource metrics~\cite{burns2016borg}. However, the reactive nature of HPA---which scales only \textit{after} a threshold is breached---creates a critical vulnerability window during sudden traffic surges~\cite{nguyen2022survey}.

This vulnerability is especially acute in latency-sensitive applications such as e-commerce platforms, where a P95 latency exceeding 500\,ms directly impacts user experience and revenue~\cite{amazon_latency}. The root cause lies in HPA's feedback loop: metrics are scraped at 15--30 second intervals, evaluated against static thresholds, and only then are scaling decisions enacted. Including pod scheduling, image pulling, and container initialization, the total delay from surge onset to pod readiness typically spans 60--90 seconds---an eternity for services operating under strict latency budgets.

\textbf{Contribution.} We propose a proactive autoscaling framework that replaces the reactive feedback loop with a predictive controller powered by Transformer-based time-series forecasting. Our key contributions are:

\begin{enumerate}
    \item A real-world experimental framework deployed on Google Kubernetes Engine (GKE) using two heterogeneous microservice benchmarks---the Google Online Boutique (10 services) and Weaveworks Sock Shop (13+ services)---with a combined dataset of over 35,000 temporally aligned observations collected over 127 hours of continuous traffic simulation.
    \item A comparative evaluation of two Transformer architectures---vanilla Transformer encoder and PatchTST~\cite{nie2023patchtst}---for CPU utilization forecasting in Kubernetes environments, demonstrating up to 98.4\% reduction in prediction MAE and 98.2\% spike detection recall.
    \item A cross-platform generalizability study demonstrating that a PatchTST model trained on one benchmark achieves a 69\% improvement over HPA on a structurally different benchmark in a zero-shot transfer setting, establishing that the learned representations capture universal microservice traffic invariants.
\end{enumerate}

The remainder of this paper is organized as follows: Section~\ref{sec:related} surveys related work. Section~\ref{sec:architecture} describes the system architecture. Section~\ref{sec:methodology} details the methodology. Section~\ref{sec:results} presents experimental results. Section~\ref{sec:discussion} discusses findings and limitations. Section~\ref{sec:conclusion} concludes.

% ==============================================================================
\section{Related Work}\label{sec:related}
% ==============================================================================

\subsection{Kubernetes Autoscaling}
The native Kubernetes HPA implements a proportional control loop that adjusts replica counts based on the ratio of current to target CPU utilization~\cite{k8s_hpa}. While effective for steady-state workloads, several studies have documented its limitations under bursty traffic patterns~\cite{nguyen2022survey, lorido2014review}. Vertical Pod Autoscaler (VPA) adjusts resource requests per pod but cannot respond to rapid demand changes~\cite{rzadca2020autopilot}. Custom metrics autoscaling extends HPA with application-level indicators but remains fundamentally reactive~\cite{baarzi2021showar}.

\subsection{Predictive Autoscaling}
Several approaches have explored predictive autoscaling using machine learning. LSTM-based predictors~\cite{imdoukh2020machine} demonstrated the feasibility of forecasting CPU usage but suffered from vanishing gradient issues over long horizons. Reinforcement learning approaches~\cite{qiu2020firm} offer adaptive policies but require extensive training environments and exhibit instability during deployment. Yan et al.~\cite{yan2021hansel} proposed a hybrid approach combining reactive scaling with Bi-LSTM forecasting, achieving moderate improvements. However, recurrent models lack the attention mechanism's ability to capture long-range temporal dependencies efficiently.

\subsection{Transformer Models for Time-Series}
The Transformer architecture~\cite{vaswani2017attention} has been widely adopted for sequence modeling through its self-attention mechanism. Its application to time-series forecasting has yielded strong results across domains~\cite{wen2023transformers}. PatchTST~\cite{nie2023patchtst} introduced a patch-based approach that segments time-series into subseries-level patches, enabling the model to capture local semantic information while maintaining computational efficiency through reduced token counts. To our knowledge, this is the first application of PatchTST to Kubernetes autoscaling, and the first to evaluate cross-platform generalizability of Transformer-based autoscalers.

% ==============================================================================
\section{System Architecture}\label{sec:architecture}
% ==============================================================================

\subsection{Cluster Configuration}
Our experimental platform consists of GKE clusters provisioned with e2-standard-4 machine types (4 vCPU, 16\,GB RAM per node) in a 3-node configuration. We deploy two microservice benchmarks to evaluate cross-platform robustness:

\begin{itemize}
    \item \textbf{Google Online Boutique}~\cite{online_boutique}: A 10-service e-commerce application implemented across Go, Python, C\#, and Java, simulating a production retail platform with services for product catalog, cart management, checkout, payment processing, and shipping.
    \item \textbf{Weaveworks Sock Shop}~\cite{sock_shop}: A 13-service microservice demonstration built on a heterogeneous stack of Java (Spring Boot), Node.js, and Go, featuring a deeper inter-service dependency graph and more diverse resource profiles than Online Boutique.
\end{itemize}

Both clusters are instrumented with Prometheus and kube-state-metrics for fine-grained observability, collecting CPU, memory, and replica count metrics at 2--15 second resolution. Locust serves as the distributed load generation framework, enabling reproducible traffic patterns across both benchmarks.

\subsection{Proactive Controller Architecture}
The AI-based autoscaler operates as an external controller that replaces the HPA decision loop with a three-stage pipeline:

\begin{enumerate}
    \item \textbf{Observe}: The controller queries Prometheus at regular intervals for the current CPU utilization, memory usage, request throughput (QPS), and P95 latency of the target frontend deployment.
    \item \textbf{Predict}: The most recent 30 observations (constituting a 60-second lookback window) are fed into the trained PatchTST model to forecast CPU utilization for the next time step.
    \item \textbf{Act}: If the predicted CPU exceeds 70\%, the controller proactively scales the deployment via the Kubernetes API, provisioning pods 45--60 seconds before the actual surge materializes.
\end{enumerate}

This architecture decouples the scaling decision from the observed metric, eliminating the fundamental latency bottleneck of reactive approaches.

% ==============================================================================
\section{Methodology}\label{sec:methodology}
% ==============================================================================

\subsection{Dataset Collection}

\subsubsection{Google Online Boutique}
We collected a high-fidelity dataset by running diverse traffic patterns against the Online Boutique frontend over 117 hours of continuous operation. Prometheus metrics were extracted via the HTTP API at 2-second resolution, yielding 32,402 raw observations. These were temporally aligned with Locust traffic logs (130,524 entries) using \texttt{merge\_asof} with a 15-second tolerance window, accounting for a 5.5-hour IST--UTC timezone offset between the Locust client and the GKE cluster. After filtering zero-QPS intervals (idle periods), the final dataset contains \textbf{28,108 observations} with four features: QPS, P95 latency, memory utilization, and CPU utilization (prediction target).

\subsubsection{Weaveworks Sock Shop}
A separate data collection campaign was conducted against the Sock Shop frontend over 10 hours at 5-second resolution, yielding \textbf{7,147 observations} after temporal alignment and IST--UTC correction. The higher sampling rate was chosen to capture the rapid inter-service dependency fluctuations characteristic of the Sock Shop's deeper service mesh.

\subsection{Traffic Patterns}
Two distinct traffic patterns were employed to ensure model generalization:

\begin{itemize}
    \item \textbf{Sine Wave}: Smooth sinusoidal oscillation between 50 and 500 concurrent users with a 10-minute period, testing gradual ramp-up and ramp-down prediction.
    \item \textbf{Spike Stress}: 8 minutes of low traffic (50 users) followed by 2-minute bursts of 1,000 concurrent users, testing sudden surge prediction capability.
\end{itemize}

\subsection{Model Architectures}

\subsubsection{Vanilla Transformer Encoder}
Our baseline deep learning model employs a single Transformer encoder block with multi-head self-attention (4 heads, key dimension 64), followed by a position-wise feed-forward network (64 units, ReLU activation). Global average pooling aggregates temporal features before a dense prediction head outputs the forecasted CPU value.

\subsubsection{PatchTST}
Following Nie et al.~\cite{nie2023patchtst}, we segment the input sequence into overlapping patches of size 6 with stride 4, yielding 7 patches per 30-step window. Each patch is linearly projected to a 128-dimensional embedding and processed through 3 stacked Transformer encoder layers with 4 attention heads each. Positional encodings are added to preserve temporal ordering. Global average pooling over the patch dimension produces the final representation for CPU prediction through a linear head.

\subsection{HPA Baseline Simulation}
To ensure a fair comparison, we simulate the reactive HPA baseline by applying a temporal shift to the actual demand signal. The shift corresponds to the realistic end-to-end HPA lag of 60 seconds, comprising metric scraping (15s), threshold evaluation (15s), and pod scheduling plus startup (30s). The HPA's ``prediction'' at time $t$ is therefore the actual demand at time $t - 60\text{s}$, representing the best-case scenario for a reactive controller operating under production conditions.

\subsection{Training Configuration}
Both models were trained with the Adam optimizer~\cite{kingma2015adam}, MSE loss, and a batch size of 64. The dataset was split 80/20 for training and testing. The vanilla Transformer used early stopping with patience of 3 epochs, while PatchTST was trained for 15 epochs. All experiments were conducted with T4 GPU acceleration. The environment variable \texttt{TF\_ENABLE\_ONEDNN\_OPTS=0} was set to prevent numerical instability from fused operations.

% ==============================================================================
\section{Experimental Results}\label{sec:results}
% ==============================================================================

\subsection{Prediction Accuracy}
Table~\ref{tab:results} presents the prediction performance of the three approaches against the realistic 60-second lag HPA baseline.

\begin{table}[htbp]
\caption{Prediction Performance on Google Online Boutique (Baseline: 60s-Lag HPA)}
\label{tab:results}
\centering
\begin{tabular}{lccc}
\toprule
\textbf{Metric} & \textbf{HPA (60s)} & \textbf{Transformer} & \textbf{PatchTST} \\
\midrule
MSE & 0.6124 & \textbf{0.00019} & 0.00028 \\
MAE & 0.8241 & 0.0155 & \textbf{0.0128} \\
RMSE & 0.7826 & 0.0138 & \textbf{0.0179} \\
\midrule
MSE Reduction & --- & 99.97\% & 99.95\% \\
MAE Reduction & --- & 98.1\% & \textbf{98.4\%} \\
\bottomrule
\end{tabular}
\end{table}

The vanilla Transformer achieves the lowest MSE (0.00019), indicating superior ability to minimize large prediction errors. PatchTST achieves the lowest MAE (0.0128), demonstrating more consistent predictions with fewer outlier deviations. Both architectures reduce the prediction error by over two orders of magnitude relative to the reactive baseline.

\subsection{Live Cluster Performance}
In live deployment tests, the proactive controller was benchmarked against standard HPA across five independent runs using the Spike Stress pattern (Table~\ref{tab:latency}).

\begin{table}[htbp]
\caption{P95 Latency Comparison ($n = 5$ independent runs)}
\label{tab:latency}
\centering
\begin{tabular}{lcc}
\toprule
\textbf{Metric} & \textbf{Standard HPA} & \textbf{AI Controller} \\
\midrule
Mean P95 Latency & 1,020 ms & \textbf{215 ms} \\
Std. Deviation ($\sigma$) & 45 ms & \textbf{12 ms} \\
Latency Reduction & --- & \textbf{78.9\%} \\
\bottomrule
\end{tabular}
\end{table}

The AI controller reduces mean P95 latency by 78.9\%, from 1,020\,ms to 215\,ms, while also exhibiting substantially lower variance ($\sigma = 12$\,ms vs.\ 45\,ms), indicating more predictable and consistent performance under load.

\subsection{Scaling Lead Time}
Analysis of scaling events reveals that the AI controller initiates pod provisioning 45--60 seconds before the traffic surge, compared to HPA which begins scaling 15--30 seconds \textit{after} the surge onset. This 60--90 second advantage ensures pods are fully initialized and ready to serve traffic at the moment demand peaks.

\subsection{Cross-Platform Robustness}
To evaluate whether the model's accuracy is specific to the Online Boutique or reflects a general capability, we trained and evaluated PatchTST independently on the Sock Shop benchmark. Table~\ref{tab:cross_app} compares the operational metrics across both platforms.

\begin{table}[htbp]
\caption{Cross-Application Operational Metrics: PatchTST vs.\ HPA}
\label{tab:cross_app}
\centering
\begin{tabular}{lcccc}
\toprule
\textbf{Metric} & \multicolumn{2}{c}{\textbf{Online Boutique}} & \multicolumn{2}{c}{\textbf{Sock Shop}} \\
\cmidrule(lr){2-3} \cmidrule(lr){4-5}
& \textbf{PatchTST} & \textbf{HPA} & \textbf{PatchTST} & \textbf{HPA} \\
\midrule
RMSE & \textbf{0.0179} & 0.7826 & \textbf{0.0917} & 1.0240 \\
Recall (Spikes) & \textbf{98.2\%} & 52.1\% & \textbf{89.5\%} & 56.8\% \\
False Positives & \textbf{1.1\%} & 48.2\% & \textbf{4.3\%} & 43.1\% \\
MAE Improvement & \multicolumn{2}{c}{97.7\%} & \multicolumn{2}{c}{85.5\%} \\
\bottomrule
\end{tabular}
\end{table}

Several observations follow from these results. First, the proactive controller achieves high spike detection recall on both benchmarks (98.2\% and 89.5\%), while HPA detects barely half of the spikes due to its inherent reaction lag. Second, the false positive rate remains below 5\% for PatchTST, compared to 43--48\% for HPA, indicating that the reactive approach maintains unnecessary over-provisioning long after traffic subsides. Third, the RMSE increase from 0.0179 (Online Boutique) to 0.0917 (Sock Shop) reflects the higher inter-service volatility in the Sock Shop's deeper dependency mesh, yet the model remains highly effective.

\subsection{Generalizability via Transfer Learning}
To assess whether the PatchTST model learns application-specific patterns or universal traffic dynamics, we conducted a zero-shot transfer learning experiment. A model trained exclusively on the 28,108-sample Online Boutique dataset was evaluated directly on the Sock Shop test set without any fine-tuning or domain adaptation.

\begin{table}[htbp]
\caption{Transfer Learning: Online Boutique $\rightarrow$ Sock Shop (Zero-Shot)}
\label{tab:transfer}
\centering
\begin{tabular}{lccc}
\toprule
\textbf{Metric} & \textbf{Source} & \textbf{Target} & \textbf{HPA} \\
& \textbf{(Boutique)} & \textbf{(Sock Shop)} & \textbf{Baseline} \\
\midrule
MSE & 0.0003 & 0.0258 & --- \\
MAE & 0.0104 & 0.1582 & 0.5101 \\
\midrule
MAE vs.\ HPA & --- & \textbf{69.0\%} & --- \\
\bottomrule
\end{tabular}
\end{table}

Despite a natural domain gap---the MAE increases from 0.0104 (in-domain) to 0.1582 (cross-domain)---the transferred model still outperforms the reactive HPA baseline by 69\%. This result provides strong evidence that PatchTST captures generalizable temporal patterns inherent to microservice workloads, rather than memorizing application-specific resource signatures. The transferred model's performance, while lower than a domain-specific model (MAE 0.074), remains operationally viable and could serve as a strong initialization for rapid fine-tuning on new deployments.

% ==============================================================================
\section{Discussion}\label{sec:discussion}
% ==============================================================================

\subsection{Why Proactive Outperforms Reactive}
The fundamental advantage of our approach lies in eliminating the HPA feedback loop delay. Standard HPA requires: (1) metric scraping (15s), (2) threshold evaluation (15s), and (3) pod startup (30--45s)---totaling 60--75 seconds of unprotected exposure during surges. Our proactive controller compresses this to near-zero by pre-positioning resources based on predicted demand. The operational consequence is dramatic: HPA misses nearly half of all traffic spikes (52\% recall), while our controller detects over 98\%.

\subsection{Transformer vs.\ PatchTST Trade-offs}
The vanilla Transformer achieves lower MSE, indicating fewer large prediction errors, while PatchTST achieves lower MAE, reflecting greater consistency across all predictions. The patch-based approach aggregates local temporal context before applying attention, which smooths out noise from Prometheus sensor jitter and ephemeral CPU fluctuations. For autoscaling decisions, where consistent directional accuracy matters more than point precision, PatchTST's stability makes it the preferred choice for production deployment.

\subsection{Implications for Production Deployment}
The transfer learning results (Table~\ref{tab:transfer}) carry significant practical implications. In production environments, obtaining a large labeled dataset for every new microservice deployment is prohibitively expensive. Our results demonstrate that a model pre-trained on a single reference application can provide immediate, meaningful improvements over HPA for new deployments, with optional fine-tuning to reach domain-specific accuracy.

\subsection{Limitations and Future Work}
\begin{itemize}
    \item While we evaluate on two benchmarks, both represent e-commerce workloads. Extending to other domains (video streaming, IoT, financial services) would strengthen the generalizability claim.
    \item The model requires periodic retraining as application behavior evolves over extended deployment periods.
    \item The current approach focuses on CPU-based scaling. Extending to multi-objective optimization incorporating latency SLOs, cost constraints, and energy efficiency is a natural direction for future work.
    \item The transfer learning experiment uses a single source--target pair. Evaluating multi-source transfer and domain adaptation techniques could further improve cross-platform performance.
\end{itemize}

% ==============================================================================
\section{Conclusion}\label{sec:conclusion}
% ==============================================================================

We presented a proactive Kubernetes autoscaling framework powered by PatchTST, a patch-based Transformer architecture for time-series forecasting. Through extensive experiments on real-world GKE clusters running two heterogeneous microservice benchmarks---Google Online Boutique and Weaveworks Sock Shop---we demonstrated the following:

\begin{enumerate}
    \item PatchTST reduces CPU prediction MAE by 98.4\% compared to the reactive HPA baseline (60-second systemic lag), achieving an RMSE of 0.0179 on Online Boutique and 0.0917 on Sock Shop.
    \item The proactive controller detects 89.5--98.2\% of traffic spikes with less than 5\% false positives, compared to HPA's 52\% recall and 43--48\% resource waste. In live deployment, P95 latency is reduced by 78.9\%.
    \item A zero-shot transfer learning experiment confirms that PatchTST captures generalizable microservice traffic patterns, achieving a 69\% improvement over HPA on an unseen benchmark without any fine-tuning.
\end{enumerate}

These results establish Transformer-based proactive autoscaling as a robust, generalizable, and operationally superior alternative to reactive HPA for latency-sensitive microservice deployments.

% ==============================================================================
% REFERENCES
% ==============================================================================
\begin{thebibliography}{00}

\bibitem{dragoni2017microservices} N. Dragoni et al., ``Microservices: Yesterday, Today, and Tomorrow,'' in \textit{Present and Ulterior Software Engineering}, Springer, 2017, pp.~195--216.

\bibitem{burns2016borg} B. Burns et al., ``Borg, Omega, and Kubernetes: Lessons learned from three container-management systems over a decade,'' \textit{ACM Queue}, vol.~14, no.~1, pp.~70--93, 2016.

\bibitem{nguyen2022survey} T.-T. Nguyen et al., ``A Survey on Autoscaling in Kubernetes,'' \textit{IEEE Access}, vol.~10, pp.~678--702, 2022.

\bibitem{amazon_latency} Amazon, ``How Loading Time Affects Your Bottom Line,'' Amazon Web Services Blog, 2019.

\bibitem{nie2023patchtst} Y. Nie et al., ``A Time Series is Worth 64 Words: Long-term Forecasting with Transformers,'' in \textit{Proc. ICLR}, 2023.

\bibitem{k8s_hpa} Kubernetes Authors, ``Horizontal Pod Autoscaler,'' Kubernetes Documentation, 2024.

\bibitem{lorido2014review} T. Lorido-Botr\'an et al., ``A Review of Auto-scaling Techniques for Elastic Applications in Cloud Environments,'' \textit{J. Grid Comput.}, vol.~12, no.~4, pp.~559--592, 2014.

\bibitem{rzadca2020autopilot} K. Rzadca et al., ``Autopilot: Workload Autoscaling at Google Scale,'' in \textit{Proc. EuroSys}, 2020, pp.~1--16.

\bibitem{baarzi2021showar} A.~F. Baarzi and G. Kesidis, ``SHOWAR: Right-Sizing and Efficient Scheduling of Microservices,'' in \textit{Proc. SoCC}, 2021.

\bibitem{imdoukh2020machine} M. Imdoukh et al., ``Machine Learning-Based Auto-Scaling for Containerized Applications,'' \textit{Neural Comput. \& Applic.}, vol.~32, pp.~9745--9760, 2020.

\bibitem{qiu2020firm} H. Qiu et al., ``FIRM: An Intelligent Fine-Grained Resource Management Framework for SLO-Oriented Microservices,'' in \textit{Proc. OSDI}, 2020, pp.~805--825.

\bibitem{yan2021hansel} M. Yan et al., ``Hansel: Adaptive Horizontal Scaling of Microservices Using Bi-LSTM,'' in \textit{Proc. ICSOC}, 2021, pp.~107--122.

\bibitem{vaswani2017attention} A. Vaswani et al., ``Attention is All You Need,'' in \textit{Proc. NeurIPS}, 2017, pp.~5998--6008.

\bibitem{wen2023transformers} Q. Wen et al., ``Transformers in Time Series: A Survey,'' in \textit{Proc. IJCAI}, 2023.

\bibitem{online_boutique} Google Cloud, ``Online Boutique: Cloud-Native Microservices Demo Application,'' GitHub, 2024.

\bibitem{sock_shop} Weaveworks, ``Sock Shop: A Microservices Demo Application,'' GitHub, 2024.

\bibitem{kingma2015adam} D. P. Kingma and J. Ba, ``Adam: A Method for Stochastic Optimization,'' in \textit{Proc. ICLR}, 2015.

\end{thebibliography}

\end{document}
